\section{Cuts in the original variables}\label{sec:original}

Proposition~\ref{prop:support} produces inequalities in lifted coordinates
$(x,w)$. Using them requires variables $w_j$ with $w_j=f_j(x)$ in the model.
If the solver's representation does not already contain exactly these
variables, adding them changes the model: new equalities, new variables for
branching and propagation, and possibly a different presolve. Our first
implementation used such a reformulation and needed a control mode to
separate its effect from that of the cuts (Appendix~\ref{app:campaigns12}).
This section shows how to avoid the auxiliary variables. The nonlinear
functions are minimized jointly over the block domain, and the rows of the
model then eliminate them. The underlying facts are classical Lagrangian
duality; we state them in the form that the implementation and its
certificates need.

\subsection{Aggregating nonlinear rows}\label{sec:aggregation}

Let $v\in\R^N$ be the vector of all model variables, including an objective
variable if there is one, and let $x=Sv$ be a block of $n$ of them, selected
by a $0/1$ matrix $S$. Suppose the model contains rows
\begin{equation}\label{eq:rows}
g(Sv)+Av\le b
\end{equation}
with $g:D\to\R^m$ and affine remainders $Av$ that may involve any variables,
including block variables; an equality is written as two opposite rows. Here
$D$ is a compact set that contains the block part $Sv$ of every point to
which the cuts must apply, for example the block box intersected with the
affine rows that involve only block variables.

\begin{proposition}[Aggregated support cut]\label{prop:aggregate}
Let $a\in\R^n$, $\lambda\in\R^m_{\ge0}$ and
\begin{equation}\label{eq:agg-support}
\beta\ \le\ \varphi(a,\lambda):=\inf_{u\in D}\bigl\{a\T u+\lambda\T g(u)\bigr\}.
\end{equation}
Then every point that satisfies~\eqref{eq:rows} and $Sv\in D$ satisfies
\begin{equation}\label{eq:agg-cut}
(S\T a-A\T\lambda)\T v\ \ge\ \beta-\lambda\T b .
\end{equation}
\end{proposition}

\begin{proof}
At such a point, $\lambda\ge0$ gives $\lambda\T g(Sv)\le\lambda\T(b-Av)$.
Substitute this into $a\T Sv+\lambda\T g(Sv)\ge\beta$.
\end{proof}

This is weak Lagrangian duality
\citep[Sections~5.1.3 and~5.2.2]{BoydVandenberghe2004}. It needs neither
convexity nor optimal multipliers. The same inequality is the Lagrangian
cut of a block in the extended block-separable reformulation of
\citet[Section~7.1]{Nowak2005}, after the block's auxiliary variables are
eliminated with the rows~\eqref{eq:rows}; related cuts and range reductions
appear in \citet{TawarmalaniSahinidis2004}, \citet{KaruppiahGrossmann2008}
and \citet{DomesNeumaier2016}. Its strength comes from minimizing the
nonlinear functions jointly over their common domain before they are
eliminated. Nonnegative combinations of linear rows that are already in the
relaxation cannot cut off any point of it.

An objective is handled through its epigraph. For minimization of
$g_0(x)+\ell_0(v)+c_0$ with objective variable $t$, the row
$g_0(x)+\ell_0(v)-t\le-c_0$ is one of the rows~\eqref{eq:rows}; for
maximization the hypograph row $-g_0(x)-\ell_0(v)+t\le c_0$ is used. The two
multipliers of the rows of an equality may be combined into one free
multiplier. Validity uses only $Sv\in D$ and the selected rows. Integrality,
the other constraints and the objective play no role, so the cut is valid
for the mixed-integer feasible set and for its closed convex hull. If $D$ is
built from bounds valid only at a node of the search tree, or from an
incumbent cutoff, the cut is valid only there; our implementation uses only
original bounds and rows, so its cuts are global.

The choice of $D$ matters: affine rows that involve only block variables
belong in $D$, not only in the LP. On the triangle $\{(x,y):x,y\ge0,\ x+y\le1\}$
every point satisfies $x^2+2xy+y^2\le x+y$ and $xy\le\tfrac14$, because
$s=x+y\in[0,1]$ gives $s^2\le s$ and $4xy\le s^2$. In the coordinates
$(m_x,m_y,s_x,s_y,p_{xy})$ these are linear inequalities valid for the hull
of the quadratic graph over the triangle. The point
$(\tfrac12,\tfrac12,\tfrac12,\tfrac12,\tfrac12)$, the average of the graph points
at $(0,0)$ and $(1,1)$ of the box $[0,1]^2$, lies in the hull over the box and
satisfies $m_x+m_y\le1$, but violates both inequalities. Intersecting the
hull over the box with the row is therefore strictly weaker than
convexifying over the row-constrained domain, as is known for factorable
relaxations \citep{Anstreicher2012,ZhuHeTawarmalani2026,WuMutsNowakHendrix2025}.
Here the first inequality is the sum of the RLT products $(1-x-y)x\ge0$ and
$(1-x-y)y\ge0$, so the example shows the weakness of the box hull, not an
advantage over RLT or semidefinite relaxations; for a single bilinear term,
SCIP's projection-based envelopes \citep{MullerSerranoGleixner2020} exploit
linear rows in the same spirit.

\subsection{What the cuts can express}\label{sec:closure}

Let $K=\conv\{(u,g(u)):u\in D\}$ and $Q=\{0\}\times\R^m_{\ge0}$. The affine
map $T(v)=(Sv,\ b-Av)$ places a model point in the lifted space of the block.
The next theorem identifies the set described by all cuts~\eqref{eq:agg-cut}.
It is the set version of the classical fact that the Lagrangian dual bound of
a nonconvex problem equals the optimal value of its convexification in the
joint space of variables and constraint values
\citep{Falk1969,Geoffrion1974,FeltenmarkKiwiel2000,LemarechalRenaud2001}; the
separation argument is the one that \citet[Theorem~3]{ChenLuedtke2022} use
for Lagrangian cuts in two-stage integer programs. We include the short proof
because our setting has nonlinear rows and affine remainders on arbitrary
variables, and because no constraint qualification is needed.

\begin{theorem}[Closure of the aggregated cuts]\label{thm:closure}
Let $D\subset\R^n$ be nonempty and compact and let $g$ be continuous on $D$.
A point $v$ satisfies every inequality~\eqref{eq:agg-cut} with
$\beta=\varphi(a,\lambda)$, over all $a\in\R^n$ and $\lambda\in\R^m_{\ge0}$, if
and only if
\begin{equation}\label{eq:closure}
v\in\mathcal C:=T^{-1}(K+Q)=\bigl\{v:\ \exists\,y\ \text{with}\ (Sv,y)\in K\ \text{and}\ y\le b-Av\bigr\}.
\end{equation}
\end{theorem}

\begin{proof}
The set $K$ is compact and convex, and $Q$ is a closed convex cone, so
$U=K+Q$ is closed and convex. A linear functional $(a,\lambda)$ has finite
infimum on $U$ only if $\lambda\ge0$, and then the infimum is
$\varphi(a,\lambda)$. Evaluating the valid inequality
$a\T u+\lambda\T y\ge\varphi(a,\lambda)$ of $U$ at $T(v)$ gives
exactly~\eqref{eq:agg-cut}, so every point of $\mathcal C$ satisfies all cuts.
Conversely, if $T(v)\notin U$, a linear functional strictly separates $T(v)$
from $U$; its infimum on $U$ is finite, so $\lambda\ge0$, and the
corresponding cut is violated at $v$.
\end{proof}

If auxiliary variables $w=g(x)$ were added, the joint-graph relaxation of the
block would be $(x,w)\in K$, $w+Av\le b$, and its projection onto $v$ is
$\mathcal C$. The direct cuts can therefore express exactly what this lifted
relaxation expresses about the original variables, without adding variables
or equalities to the model, and no more. Relaxations that keep the rows
inside the block, such as the convex-hull relaxations of block-decomposition
methods \citep[Section~3]{Nowak2005}, \citep{WuMutsNowakHendrix2025}, can be
strictly stronger (Example~\ref{ex:quarter}). There is also a cost: with the
same linearization points, outer approximation of a lifted formulation can be
tighter than outer approximation in the original space
\citep{TawarmalaniSahinidis2005}, and a finite set of direct cuts may need
more members than a lifted description. For one row, $\mathcal C$ is the
familiar envelope relaxation. Write $\operatorname{vex}g$ and
$\operatorname{cav}g$ for the convex and concave envelopes of a scalar
function $g$ over $D$, so that $\operatorname{vex}g(x)=\min\{y:(x,y)\in K\}$.

\begin{corollary}[One row]\label{cor:one-row}
If $m=1$, then $\mathcal C=\{v:\ Sv\in\conv D,\ \operatorname{vex}g(Sv)+Av\le b\}$.
For an equality $g(Sv)+Av=b$, written as two rows, $\mathcal C$ is the set of
$v$ with $Sv\in\conv D$ and
$\operatorname{vex}g(Sv)\le b-Av\le\operatorname{cav}g(Sv)$.
\end{corollary}

\begin{proof}
The projection of $K$ onto $x$ is $\conv D$, and for each $x\in\conv D$ the
fiber $\{y:(x,y)\in K\}$ is the interval
$[\operatorname{vex}g(x),\operatorname{cav}g(x)]$. A suitable $y$
in~\eqref{eq:closure} exists if and only if $\operatorname{vex}g(Sv)\le b-Av$.
For the equality, the lifted point of the two rows is $(y,-y)$ with $y$ in the
same fiber, and the two conditions $y\le b-Av$ and $-y\le-(b-Av)$ force
$y=b-Av$.
\end{proof}

Joint treatment can only help when rows share block variables or when $D$
couples them: if $D$ is a product $D^1\times D^2$ and the rows split into two
groups whose nonlinear functions depend only on $x^1$ and only on $x^2$, then
$\mathcal C$ is the intersection of the closures of the two groups, because the
convex hull of a product of sets is the product of their convex hulls.
When the rows do share variables, the gain can be strict.

\begin{example}[Two rows sharing one variable]\label{ex:two-rows}
Let $D=[-1,1]$ and consider the rows $x^2-v_1\le0$ and $-x^3+v_2\le0$. The
convex envelope of $x^2$ is $x^2$, and the convex envelope of $-x^3$ on
$[-1,1]$ takes the value $-1/4$ at $x=0$ (on $[-1/2,1]$ it is the line through
$(1,-1)$ that is tangent to $-x^3$ at $x=-1/2$). By
Corollary~\ref{cor:one-row}, the intersection of the two one-row closures
contains $(x,v_1,v_2)=(0,0,1/4)$. The aggregated cut with $a=0$ and
$\lambda=(1,1)$ uses $\min_{u\in[-1,1]}(u^2-u^3)=0$ and reads $v_1-v_2\ge0$,
which this point violates. In terms of~\eqref{eq:closure}, the only
probability measure on $[-1,1]$ with mean $0$ and second moment $0$ is the
point mass at $0$, whose third moment is $0$, not $1/4$.
\end{example}

That joint relaxations dominate separate envelopes is well known
\citep[Lemma~3.4]{Nowak2005}, \citep{Tawarmalani2010,LiersEtAl2021}; the
example makes the effect concrete for the cut family used here.

\subsection{Relation to the hull of the feasible set}\label{sec:feasible-hull}

The closure $\mathcal C$ is a relaxation of the set
$\Sigma=\{v:\ Sv\in D,\ v\text{ satisfies~\eqref{eq:rows}}\}$, so
$\conv\Sigma\subseteq\mathcal C$. Equality holds when every row has its own
free direction in the remainders.

\begin{proposition}[Free remainders]\label{prop:free-remainders}
Write $v=(x,z)$ with $x$ the block, and let $M$ be the matrix formed by the
columns of $A$ that belong to $z$, with one row for each inequality and one
row for each equality (of the two opposite rows of an equality, $M$ keeps one).
If $M$ has full row rank, then $\mathcal C=\conv\Sigma$.
\end{proposition}

\begin{proof}
Let $v=(x,z)\in\mathcal C$. By Carath\'eodory's theorem there are $u_j\in D$ and
weights $\theta_j>0$ summing to one with $\sum_j\theta_ju_j=x$ and
$\bar y:=\sum_j\theta_jg(u_j)\le b-Av$, with equality in the rows of each
equality. Let $A_x$ denote the block columns of $A$, restricted like $M$, let
$M^+$ be a right inverse of $M$, put $\rho_j=\bar y-g(u_j)+A_x(x-u_j)$ in the
same rows, and $z_j=z+M^+\rho_j$. A direct substitution shows that each row
of $(u_j,z_j)$ has the same left-hand side as $\bar y+Av$, so
$(u_j,z_j)\in\Sigma$. Since $\sum_j\theta_j\rho_j=0$, the convex combination of
the points $(u_j,z_j)$ is $v$.
\end{proof}

The objective row always contains $-t$, so for the objective alone the
closure is the hull of the epigraph over $D$. The setting of
\citet{LiersEtAl2021}, rows $z_j=g_j(x)$ with a separate variable $z_j$ for
each row, is also covered by Proposition~\ref{prop:free-remainders}. Without
free remainders the gap can be strict.

\begin{example}[Gap to the feasible hull]\label{ex:quarter}
Let $D=[0,1]$ and take the equality row $x^2=1/4$, whose feasible set is
$\{1/2\}$. By Corollary~\ref{cor:one-row}, $\mathcal C=\{x:x^2\le1/4\le x\}=[1/4,1/2]$,
and the two cuts $x\ge1/4$ and $x\le1/2$ describe it. The point $x=1/4$ is the
mean of the measure with mass $3/4$ at $0$ and $1/4$ at $1$, which satisfies
the row in expectation.
\end{example}

This is a Lagrangian duality gap \citep{Geoffrion1974,Nowak2005}: the rows are
imposed on averages, after convexification. Even the complete family of cuts
accepts $x=1/4$, so the gap is not a failure of direction search. It is
closed by convexifying aggregated sets instead of aggregated functions, as in
the aggregation closure of \citet{DeyMunozSerrano2022}: here
$\conv\{x\in[0,1]:x^2\le1/4\}=[0,1/2]$ and $\conv\{x\in[0,1]:x^2\ge1/4\}=[1/2,1]$
intersect in $\{1/2\}$. It is also closed by placing the row inside $D$, which
changes the support problem; our implementation keeps nonlinear rows out of
$D$ so that the support problems stay in the classes of
Section~\ref{sec:quadratic}. With
$\Sigma_\lambda=\{v:\ Sv\in D,\ \lambda\T(g(Sv)+Av-b)\le0\}$ for $\lambda\ge0$, the
hierarchy is
\[
\conv\Sigma\ \subseteq\ \bigcap_{\lambda\ge0}\conv\Sigma_\lambda\ \subseteq\ \mathcal C .
\]
The middle set is the convexified closure of the surrogate relaxations of
\citet{MullerEtAl2022}, and the gap to $\mathcal C$ is a surrogate--Lagrangian
gap. Example~\ref{ex:quarter} shows that the second inclusion can be strict.
The first can be strict as well, even when $D$ is exactly the projection of
$\Sigma$: with $D=[0,2]$, a free $z$ and the rows $z\le(x-2)^2/2$ and
$z\le x^2/2$, the set $\Sigma$ satisfies $z\le\min\{x/2,1-x/2\}$ on its hull,
so $z\le1/2$ at $x=1$. Every aggregated row reads $z\le q_\lambda(x)$ with
$q_\lambda=[\lambda_1(x-2)^2+\lambda_2x^2]/(2(\lambda_1+\lambda_2))$ convex and
$q_\lambda(0)+q_\lambda(2)=2$, so the chord gives $(1,1)\in\conv\Sigma_\lambda$ for
every $\lambda$, and $(1,1)$, the mean of the point masses at $x=0$ and $x=2$,
also lies in $\mathcal C$.
